\section{优化方法}

第 2 章已经将多机械臂分类搬运问题建模为带有可选执行模式、并行资源互斥、序列相关 setup、中心安全区互斥以及序贯目标的调度优化模型。本章在该模型基础上说明具体求解方法。整体思路分为两层：第一层求解离散调度问题，确定每个任务由哪些机械臂执行以及各机械臂的作业顺序；第二层在离散调度结果给定后，进一步分析二臂/三臂系统选择和速度--能耗连续优化问题。

\subsection{总体求解框架}

本文的优化流程如图 \ref{fig:method_framework} 所示。首先根据随机生成的方块位置、方块类型和机械臂数量构造任务实例；随后为每个任务生成可行执行模式，包括单臂执行模式和双臂协同执行模式；在此基础上分别采用基础整数约束规划、自编分枝定界算法和启发式方法求解调度问题；最后利用批量算例结果讨论二臂/三臂模式选择，并在给定调度结构上进行速度--能耗非线性优化。

\begin{figure}[H]
    \centering
    \fbox{
    \begin{minipage}{0.88\textwidth}
        \centering
        任务实例生成：方块类型、初始位置、目标收集盒、机械臂布局

        $\Downarrow$

        可行模式构造：机械臂组合、处理时间 $t_k$、处理能耗 $e_k$、中心区标记 $z_k$

        $\Downarrow$

        离散调度求解：CP-SAT 基准模型 / 自编分枝定界 / 启发式与 warm start

        $\Downarrow$

        结果统计比较：$C_{\max}$、$E_{\mathrm{tot}}$、负载均衡、计算时间

        $\Downarrow$

        扩展分析：二臂/三臂模式选择、速度--能耗非线性优化、鲁棒速度选择
    \end{minipage}
    }
    \caption{本文优化方法总体框架}
    \label{fig:method_framework}
\end{figure}

上述流程体现了本文方法的分工。CP-SAT 模型用于提供可靠的基准结果；自编分枝定界算法用于展示对问题结构的显式利用，并验证不依赖通用求解器时仍可得到最优调度；启发式方法用于较大规模算例中快速获得高质量可行解，同时可作为精确搜索的初始上界。二臂/三臂模式选择和速度--能耗优化不改变第 2 章的基本调度约束，而是在调度结果之上进一步回答“是否值得增加一只机械臂”和“在给定时间要求下如何选择运行速度”的问题。

\subsection{基础方法：整数约束规划与序贯目标规划}

基础方法直接求解第 2 章建立的整数调度模型。由于该问题同时包含 0-1 模式选择、时间区间安排、机械臂资源互斥、中心安全区互斥和序列相关 setup，本文采用 OR-Tools CP-SAT 作为基础求解器。CP-SAT 适合处理此类混合了逻辑约束和整数时间变量的离散优化问题，因此可作为后续自编算法和启发式算法的对照基准。

\subsubsection{可选区间与资源互斥}

对每个可行执行模式 $k$，基础模型建立一个选择变量 $x_k$ 和一个可选时间区间 $[S_k,E_k)$。若 $x_k=1$，该区间表示任务真实执行的处理段；若 $x_k=0$，该模式不参与调度。每个任务必须且只能选择一个模式：
\[
    \sum_{k\in\mathcal{M}_j}x_k=1,\quad \forall j\in\mathcal{J}.
\]

机械臂资源约束通过 NoOverlap 实现。若模式 $k$ 占用机械臂 $a$，则其处理区间被加入机械臂 $a$ 的资源列表，同一机械臂上的任务区间不能重叠。中心安全区也采用类似的互斥资源思想：当模式 $k$ 的搬运路径经过中心安全区时，将其任务处理区间加入中心区资源列表，从而保证不同任务组不能同时占用中心区。对于同一个双臂协同任务，参与的两只机械臂属于同一任务组，不构成彼此冲突。

\subsubsection{序列相关 setup}

任务之间的空载转移取决于同一机械臂上的相邻任务顺序，因此模型需要同时决定“做哪些任务”和“按什么顺序做”。若机械臂 $a$ 在模式 $u$ 后紧接执行模式 $v$，则加入顺序约束
\[
    S_v\ge E_u+\tau_{u,v},
\]
其中 $\tau_{u,v}$ 为从任务 $u$ 的目标收集盒移动到任务 $v$ 初始位置的空载转移时间。相应的空载转移能耗 $\epsilon_{u,v}$ 也计入系统总能耗。若某任务是机械臂 $a$ 的首个任务，则 setup 从该机械臂待机位置到方块初始位置计算。

\subsubsection{三阶段序贯目标}

本文关注最大完工时间、总能耗和机械臂负载均衡三个指标。若直接加权求和，权重选择会影响结果解释；因此基础方法采用序贯目标规划，按照优先级分三阶段求解。

第一阶段优先优化最大完工时间：
\[
    P_1:\min \left(Md_1^+ + C_{\max}\right).
\]
第二阶段在保持第一阶段最优时间结果不变的基础上优化总能耗：
\[
    P_2:\min \left(Md_2^+ + E_{\mathrm{tot}}\right).
\]
第三阶段在不破坏前两级主要目标的基础上优化负载均衡：
\[
    P_3:\min \left(Md_3^+ + B\right).
\]

这种顺序体现了本文的调度偏好：首先保证任务整体尽快完成，其次在不牺牲完工时间的前提下降低能耗，最后再改善不同机械臂之间的利用均衡。基础方法的结果可作为小规模和中等规模算例中的参考解。

\subsection{自编精确方法：隐枚举与分枝定界}

为了体现对模型结构的直接利用，本文进一步实现了自编精确搜索算法。该算法不调用通用整数规划求解器，而是在任务模式和任务顺序空间中进行隐枚举，并通过分枝定界和状态支配剪枝减少搜索量。搜索目标采用与序贯目标一致的字典序：
\[
    (C_{\max},E_{\mathrm{tot}},B).
\]
即先比较最大完工时间；若完工时间相同，再比较总能耗；若二者仍相同，最后比较负载不均衡度。

\subsubsection{搜索状态与扩展规则}

一个部分调度状态记录已经安排的任务集合、每只机械臂的当前可用时间、每只机械臂的当前位置、当前累计能耗、当前最大完工时间、中心区已占用时间段以及已经生成的调度序列。

从某个部分状态出发，算法枚举尚未安排任务的可行执行模式。对于单臂任务，任务开始时间由对应机械臂完成前置 setup 后的最早可用时间决定；对于双臂协同任务，任务必须等待参与的两只机械臂均完成前置 setup 后才能同步开始。若该任务经过中心安全区，还需要将任务处理区间推迟到不与已有中心区占用区间重叠的位置。

上述过程可概括为以下步骤：
\begin{enumerate}
    \item 从空调度状态出发，初始化各机械臂可用时间和当前位置；
    \item 选择一个未安排任务，枚举其所有可行执行模式；
    \item 计算该模式在当前状态下的 setup 时间、开始时间、结束时间和增量能耗；
    \item 若涉及中心安全区，则调整开始时间以满足中心区互斥；
    \item 更新机械臂状态、中心区占用记录和当前调度序列；
    \item 利用下界和支配关系判断是否继续扩展该分支。
\end{enumerate}

\subsubsection{下界剪枝}

分枝定界的关键是为未完成的部分调度构造安全下界。安全下界表示真实最优扩展不可能优于该估计，因此若下界已经不优于当前最好完整解，该分支即可被剪去。本文采用的下界包括：
\begin{itemize}
    \item 剩余任务最早单独完成时间下界；
    \item 剩余处理工作量相对于机械臂数量的资源容量下界；
    \item 某些任务必须占用特定机械臂时的强制机械臂下界；
    \item 必须经过中心安全区任务的中心区串行容量下界；
    \item 剩余任务最小处理能耗下界。
\end{itemize}

这些下界只使用必然成立的信息，不人为假设未来任务的具体顺序，因此不会剪掉潜在最优解。剪枝判断同样采用字典序：若当前状态的下界三元组不优于已有最好解，则停止扩展该状态。

\subsubsection{状态支配剪枝}

除下界剪枝外，算法还利用状态支配关系减少重复搜索。若两个部分调度已经完成的任务集合相同，各机械臂当前位置相同，并且其中一个状态在机械臂可用时间、机械臂累计负载、累计能耗、当前最大完工时间和中心区释放时间等指标上均不劣于另一个状态，则较差状态在后续扩展中不可能得到更优结果，可以被支配剪去。

状态支配只在“已完成任务集合”和“机械臂当前位置”一致时使用，因此比较的是具有相同未来可选任务和相同出发位置的状态。这一条件保证了支配剪枝的保守性。

\subsection{启发式方法与 warm start}

精确搜索能够证明最优性，但当任务数增加、三臂系统可选组合增多时，搜索空间会迅速扩大。为提高较大规模算例的求解效率，本文实现了大邻域搜索启发式方法，并将其结果用于精确搜索的初始上界。

\subsubsection{大邻域搜索启发式}

启发式方法首先构造多个初始可行解。构造策略包括优先缩短完工时间、优先减少中心区冲突、优先降低能耗和优先改善负载均衡等。随后采用大邻域搜索思想：从当前完整调度中移除一部分任务，再将这些任务按较优位置重新插入，形成新的完整调度。

若新调度在字典序目标上优于当前解，则接受该解；在搜索前期也允许接受少量非劣或略差解，以增加跳出局部最优的机会。搜索后期则转向更保守的改进规则，并在不增加 $C_{\max}$ 的前提下尝试进一步降低能耗。该方法不证明全局最优性，但能够快速给出可行解，适合用于较大算例的对比实验。

\subsubsection{warm start 初始上界}

warm start 的思想是先用启发式方法得到一个较好的完整调度，将其作为自编分枝定界算法的初始 incumbent。这样精确搜索一开始就拥有较强上界，许多下界已经不可能优于该上界的分支可以被更早剪去。

需要强调的是，启发式解只作为初始上界使用，不改变精确算法的可行域、下界定义和支配剪枝规则。因此，只要分枝定界搜索最终完成，所得结果仍然是由精确搜索证明的最优解；warm start 只影响搜索效率，不影响最优性结论。

\subsection{二臂/三臂模式选择方法}

完成调度求解后，本文进一步比较两臂系统和三臂系统在相同任务结构下的表现。对于同一组方块数量和随机 seed，分别运行二臂调度和三臂调度，记录最大完工时间 $C_{\max}$ 和系统总能耗 $E_{\mathrm{tot}}$。再对同一任务规模下多个随机 seed 的结果取平均，以降低单个随机实例对结论的影响。

为了刻画三臂系统相对于两臂系统的时间收益，定义时间节省率：
\[
    \eta_T=\frac{C_{\max}^{2B}-C_{\max}^{3B}}{C_{\max}^{2B}}\times 100\%.
\]
若 $\eta_T$ 越大，说明增加第三只机械臂对缩短完工时间越有利。

同时定义能耗增加率：
\[
    \eta_E=\frac{E_{\mathrm{tot}}^{3B}-E_{\mathrm{tot}}^{2B}}{E_{\mathrm{tot}}^{2B}}\times 100\%.
\]
若 $\eta_E$ 越大，说明三臂系统带来的额外能耗越明显。

在不同决策偏好下，引入综合评分
\[
    S_\lambda=\eta_T-\lambda\eta_E,
\]
其中 $\lambda$ 表示对能耗代价的重视程度。当 $S_\lambda>0$ 时，说明三臂系统的时间收益足以抵消能耗代价，推荐三臂模式；当 $S_\lambda<0$ 时，说明能耗代价相对更高，推荐两臂模式。通过改变 $\lambda$，可以分析不同偏好下的系统配置选择。

\subsection{速度--能耗非线性优化}

前述调度模型中的时间和能耗参数由统一速度假设得到，主要用于解决任务分配和作业顺序问题。为了进一步分析运行速度对能耗的影响，本文在调度结构给定后构造速度--能耗非线性优化模型。该模型不重新排序任务，而是调整单臂任务和双臂协同任务的速度倍率，从而研究在截止时间约束下的经济运行速度。

\subsubsection{确定性速度优化模型}

设
\[
    s_{\mathrm{single}},\quad s_{\mathrm{dual}}
\]
分别为单臂任务和双臂协同任务的速度倍率。速度提高会缩短任务处理时间，但可能增加高速运行造成的动态代价；速度过低则会增加与任务持续时间相关的保持代价。因此本文采用 U 形速度--能耗函数。对任务类别 $r\in\{\mathrm{single},\mathrm{dual}\}$，定义
\[
    \mathcal{E}_r(s_r)=E_{r0}\left[
    \frac{\lambda_r}{s_r}
    +(1-\lambda_r)\left((1-\rho_r)+\rho_r s_r^2\right)
    \right],
\]
其中 $E_{r0}$ 为该类任务的基准能耗，$\lambda_r$ 表示低速持续时间代价权重，$\rho_r$ 表示高速动态代价权重。

给定截止时间 $D$，确定性速度优化模型可写为
\[
    \min\quad E_{\mathrm{fixed}}
    +\mathcal{E}_{\mathrm{single}}(s_{\mathrm{single}})
    +\mathcal{E}_{\mathrm{dual}}(s_{\mathrm{dual}})
\]
\[
    \text{s.t.}\quad
    C_{\mathrm{fixed}}
    +\frac{C_{\mathrm{single}}}{s_{\mathrm{single}}}
    +\frac{C_{\mathrm{dual}}}{s_{\mathrm{dual}}}
    \le D,
\]
\[
    s_{\min}\le s_{\mathrm{single}}\le s_{\max},\quad
    s_{\min}\le s_{\mathrm{dual}}\le s_{\max}.
\]

其中 $C_{\mathrm{single}}$ 和 $C_{\mathrm{dual}}$ 分别表示调度结果中可由速度倍率缩放的单臂和双臂任务时间部分，$C_{\mathrm{fixed}}$ 表示不随速度倍率变化的固定时间项；$E_{\mathrm{fixed}}$ 表示不参与速度优化的固定能耗项。

求解时，本文采用内点障碍函数处理不等式约束，并在每个障碍参数下使用最速下降方向迭代。步长通过 0.618 黄金分割一维搜索确定。得到速度解后，再通过 KKT 条件检验原始可行性、互补松弛和驻点残差，以判断数值解是否满足非线性规划的一阶最优性条件。

\subsubsection{鲁棒速度选择模型}

确定性模型通常先对多个 seed 的调度结果取平均，再进行速度优化。为了考虑随机实例波动，本文进一步构造鲁棒速度选择模型：对同一任务规模下的多个 seed 使用同一组速度倍率，并最小化最坏情况下的能耗上界。

令 $\mathcal{S}$ 表示 seed 集合，$z$ 表示最坏能耗上界，鲁棒模型为
\[
    \min z
\]
\[
    \text{s.t.}\quad
    E_s(s_{\mathrm{single}},s_{\mathrm{dual}})\le z,\quad \forall s\in\mathcal{S},
\]
\[
    C_s(s_{\mathrm{single}},s_{\mathrm{dual}})\le D,\quad \forall s\in\mathcal{S},
\]
\[
    s_{\min}\le s_{\mathrm{single}},s_{\mathrm{dual}}\le s_{\max}.
\]

该模型的含义是：面对方块随机位置带来的实例差异，选择一组统一速度，使所有样本都能在截止时间内完成，并使最坏情况下的能耗尽可能低。求解时采用逐次线性规划方法，在当前速度点对非线性约束进行一阶线性化，求解线性规划子问题，再用原非线性约束检查候选解的可行性和改进性。最终同样输出 KKT 检验结果，用于评估鲁棒速度解的可靠性。
